% ============================================================
%  The Geometry of Closure
%  Admissibility, Obstruction, and Ontological Enlargement
%  Author: Flyxion / Independent Researcher
%  Compiler: LuaLaTeX
%  Bibliography: embedded via filecontents
% ============================================================

\documentclass[12pt,oneside]{book}

% ── Fonts ────────────────────────────────────────────────────
\usepackage{fontspec}
\setmainfont{TeX Gyre Pagella}
\usepackage{microtype}

% ── Mathematics ───────────────────────────────────────────────
\usepackage{amsmath,amsthm,amssymb,mathtools}
\usepackage{unicode-math}
\setmathfont{TeX Gyre Pagella Math}

% ── Page geometry ─────────────────────────────────────────────
\usepackage[
  paperwidth=6in,paperheight=9in,
  top=1in,bottom=1in,inner=0.9in,outer=0.75in
]{geometry}
\usepackage{setspace}
\onehalfspacing

% ── Headers ───────────────────────────────────────────────────
\usepackage{fancyhdr}
\pagestyle{fancy}\fancyhf{}
\fancyhead[LE]{\small\itshape The Geometry of Closure}
\fancyhead[RO]{\small\itshape\leftmark}
\fancyfoot[C]{\small\thepage}
\renewcommand{\headrulewidth}{0.4pt}

% ── Chapter titles ────────────────────────────────────────────
\usepackage{titlesec}
\titleformat{\chapter}[display]
  {\normalfont\Large\bfseries}{\chaptertitlename\ \thechapter}{14pt}{\LARGE}
\titlespacing*{\chapter}{0pt}{28pt}{18pt}

% ── Theorem environments ──────────────────────────────────────
\theoremstyle{plain}
\newtheorem{theorem}{Theorem}[chapter]
\newtheorem{proposition}[theorem]{Proposition}
\newtheorem{lemma}[theorem]{Lemma}
\newtheorem{corollary}[theorem]{Corollary}
\newtheorem{conjecture}[theorem]{Conjecture}
\newtheorem{metatheorem}[theorem]{Meta-Theorem}

\theoremstyle{definition}
\newtheorem{definition}[theorem]{Definition}
\newtheorem{example}[theorem]{Example}
\newtheorem{construction}[theorem]{Construction}

\theoremstyle{remark}
\newtheorem{remark}[theorem]{Remark}
\newtheorem{notation}[theorem]{Notation}

% ── Hyperlinks ────────────────────────────────────────────────
\usepackage[colorlinks,linkcolor=black,citecolor=black,urlcolor=black]{hyperref}

% ── Bibliography — embedded, no external .bib needed ─────────
\usepackage{filecontents}
\begin{filecontents*}{\jobname.bib}
@book{federer1969,
  author    = {Federer, Herbert},
  title     = {Geometric Measure Theory},
  publisher = {Springer},
  address   = {New York},
  year      = {1969}
}
@article{simons1968,
  author  = {Simons, James},
  title   = {Minimal varieties in {R}iemannian manifolds},
  journal = {Annals of Mathematics},
  volume  = {88},
  number  = {1},
  pages   = {62--105},
  year    = {1968}
}
@article{bombieri1969,
  author  = {Bombieri, Enrico and De~Giorgi, Ennio and Giusti, Enrico},
  title   = {Minimal cones and the {B}ernstein problem},
  journal = {Inventiones Mathematicae},
  volume  = {7},
  pages   = {243--268},
  year    = {1969}
}
@article{bernstein1915,
  author  = {Bernstein, Serge},
  title   = {{Sur un th\'eor\`eme de g\'eom\'etrie et ses applications
             aux \'equations aux d\'eriv\'ees partielles du type elliptique}},
  journal = {Communications de la Soci\'et\'e Math\'ematique de Kharkov},
  volume  = {15},
  pages   = {38--45},
  year    = {1915}
}
@article{degiorgi1965,
  author  = {De~Giorgi, Ennio},
  title   = {Una estensione del teorema di {B}ernstein},
  journal = {Annali della Scuola Normale Superiore di Pisa},
  volume  = {19},
  pages   = {79--85},
  year    = {1965}
}
@article{almgren1966,
  author  = {Almgren, Frederick~J.},
  title   = {Some interior regularity theorems for minimal surfaces and
             an extension of {B}ernstein's theorem},
  journal = {Annals of Mathematics},
  volume  = {84},
  pages   = {277--292},
  year    = {1966}
}
@book{simon1983,
  author    = {Simon, Leon},
  title     = {Lectures on Geometric Measure Theory},
  publisher = {Australian National University},
  address   = {Canberra},
  year      = {1983}
}
@book{morgan2016,
  author    = {Morgan, Frank},
  title     = {Geometric Measure Theory: {A} Beginner's Guide},
  publisher = {Academic Press},
  edition   = {5},
  year      = {2016}
}
@article{petz1988,
  author  = {Petz, D{\'e}nes},
  title   = {Sufficiency of channels over von {N}eumann algebras},
  journal = {Quarterly Journal of Mathematics},
  volume  = {39},
  number  = {1},
  pages   = {97--108},
  year    = {1988}
}
@article{petz1986,
  author  = {Petz, D{\'e}nes},
  title   = {Quasi-entropies for finite quantum systems},
  journal = {Reports on Mathematical Physics},
  volume  = {23},
  number  = {1},
  pages   = {57--65},
  year    = {1986}
}
@book{amari2000,
  author    = {Amari, Shun-ichi and Nagaoka, Hiroshi},
  title     = {Methods of Information Geometry},
  publisher = {American Mathematical Society},
  address   = {Providence},
  year      = {2000}
}
@book{villani2009,
  author    = {Villani, C{\'e}dric},
  title     = {Optimal Transport: Old and New},
  publisher = {Springer},
  address   = {Berlin},
  year      = {2009}
}
@book{nielsen2010,
  author    = {Nielsen, Michael~A. and Chuang, Isaac~L.},
  title     = {Quantum Computation and Quantum Information},
  publisher = {Cambridge University Press},
  year      = {2010}
}
@book{maclane1998,
  author    = {Mac~Lane, Saunders},
  title     = {Categories for the Working Mathematician},
  publisher = {Springer},
  edition   = {2},
  year      = {1998}
}
@book{awodey2010,
  author    = {Awodey, Steve},
  title     = {Category Theory},
  publisher = {Oxford University Press},
  edition   = {2},
  year      = {2010}
}
@article{sinkhorn1967,
  author  = {Sinkhorn, Richard},
  title   = {Diagonal equivalence to matrices with prescribed row and
             column sums},
  journal = {American Mathematical Monthly},
  volume  = {74},
  number  = {4},
  pages   = {402--405},
  year    = {1967}
}
@article{wilde2015,
  author  = {Wilde, Mark~M. and Winter, Andreas and Yang, Dong},
  title   = {Strong converse for the classical capacity of
             entanglement-breaking and {H}adamard channels via a
             sandwiched {R}{\'e}nyi relative entropy},
  journal = {Communications in Mathematical Physics},
  volume  = {331},
  pages   = {593--622},
  year    = {2015}
}
@book{kuhn1962,
  author    = {Kuhn, Thomas~S.},
  title     = {The Structure of Scientific Revolutions},
  publisher = {University of Chicago Press},
  year      = {1962}
}
@book{lakatos1976,
  author    = {Lakatos, Imre},
  title     = {Proofs and Refutations},
  publisher = {Cambridge University Press},
  year      = {1976}
}
@article{chentsov1982,
  author  = {Chentsov, Nikolai~N.},
  title   = {Statistical Decision Rules and Optimal Inference},
  journal = {Translations of Mathematical Monographs},
  volume  = {53},
  publisher = {American Mathematical Society},
  year    = {1982}
}
@article{schrodinger1932,
  author  = {Schr{\"o}dinger, Erwin},
  title   = {{Sur la th\'eorie relativiste de l\'electron et
             l\'interpr\'etation de la m\'ecanique quantique}},
  journal = {Annales de l'Institut Henri Poincar\'e},
  volume  = {2},
  pages   = {269--310},
  year    = {1932}
}
\end{filecontents*}

\usepackage[backend=biber,style=alphabetic,sorting=nyt]{biblatex}
\addbibresource{\jobname.bib}

% ── Misc packages ─────────────────────────────────────────────
\usepackage{graphicx,xcolor,epigraph,enumitem,booktabs,array}
\setlength{\epigraphwidth}{0.64\textwidth}
\setlist[itemize]{nosep,leftmargin=1.5em}
\setlist[enumerate]{nosep,leftmargin=1.5em}

% ── Notation macros ───────────────────────────────────────────
\newcommand{\Oold}{\mathcal{O}}
\newcommand{\Onew}{\mathcal{O}'}
\newcommand{\F}{\mathcal{F}}
\newcommand{\Admiss}{\mathcal{A}}
\newcommand{\Inadmiss}{\mathcal{I}}
\newcommand{\W}{\mathcal{W}}
\newcommand{\cl}{\operatorname{cl}}
\newcommand{\dist}{\operatorname{dist}}
\newcommand{\supp}{\operatorname{supp}}
\newcommand{\tr}{\operatorname{tr}}
\newcommand{\id}{\mathrm{id}}
\newcommand{\KL}{D_{\mathrm{KL}}}
\newcommand{\Real}{\mathbb{R}}
\newcommand{\Hilb}{\mathcal{H}}
\newcommand{\density}{\mathcal{D}}
\newcommand{\bO}{\mathbf{O}}
\newcommand{\bOn}{\mathbf{O}'}

% ── Interlude environment ─────────────────────────────────────
\newenvironment{interlude}[1]{%
  \vspace{1.4em}
  \begin{center}
  \rule{0.38\textwidth}{0.4pt}\quad{\small\itshape #1}\quad\rule{0.38\textwidth}{0.4pt}
  \end{center}
  \vspace{0.4em}
  \small\itshape
}{%
  \normalsize\normalfont
  \vspace{1.4em}
}

% =============================================================
\begin{document}
\frontmatter

% ── Title page ────────────────────────────────────────────────
\begin{titlepage}
\centering\vspace*{2.2cm}
{\Huge\bfseries The Geometry of Closure}\\[1.1em]
{\Large\itshape Admissibility, Obstruction,\\[0.3em]and Ontological Enlargement}\\[3.2em]
{\large Flyxion}\\[0.5em]
{\normalsize Independent Researcher}\\[4em]
{\normalsize \today}
\vfill
{\small\itshape Working draft.}
\end{titlepage}

% ── Preface ───────────────────────────────────────────────────
\chapter*{Preface}
\addcontentsline{toc}{chapter}{Preface}

This book did not begin with geometric measure theory or quantum Bayesian
inference. It began with a pattern.

Across a sequence of projects developed under the names RSVP, Spherepop,
CLIO, TARTAN, Simulated Agency, and the Projection Crisis essays, the same
structural failure kept appearing in different contexts and under different
local names. A representation would prove locally adequate, generating reliable
predictions and useful intuitions within a familiar regime. Then something
would appear at the boundary of that regime which the representation could not
absorb. The representation did not become false. It became insufficient.
The map survived. Recovery failed.

The framework developed in this book is an attempt to say precisely what
that means. The central concepts are closure, recovery, obstruction, and
ontological commitment. They crystallized from repeated independent encounters
with the same failure mode across domains that have no obvious connection
to one another. That genealogy matters: a framework designed to fit selected
examples is always vulnerable to the charge of engineering. A framework that
emerged elsewhere and was then found to apply has different evidential status.

Two terminological notes. \emph{Collapse} is used throughout in a technical
sense: the loss of recovery while embedding is retained. The old theory
survives; it becomes a regime. \emph{Obstruction} is used in the sense
already standard in topology, category theory, and PDE: an object that
certifies, by its existence, that a certain recovery is impossible.

Three philosophical themes recur across the technical chapters, and the
reader may find it useful to name them at the outset.

First: continuation is the primitive concept, not objects. The central
question of the framework is never what a system \emph{is} but whether an
admissible continuation of it \emph{exists}. Collapse is defined by the
disappearance of admissible continuation; enlargement is defined by its
restoration.

Second: ontologies are residues, not foundations. The objects of a theory
appear fundamental only because they survive the dynamics used to generate
them. Space, probability distributions, semantic scopes, and mathematical
objects themselves are persistence structures: what remains invariant under
repeated application of governing operations.

Third: structure without blueprint. Area-minimizing currents do not know
the surface they are attempting to become. Bayesian posteriors do not know
the parameter they will converge toward. Scope graphs do not contain their
future resolutions. Global structure emerges from local admissibility
conditions, not from pre-specified endpoints. Ontological enlargement is
not the discovery of a hidden plan. It is the discovery that the admissibility
conditions were richer than previously recognized.

\vspace{1.2em}
\noindent\textit{Flyxion} \hfill Independent Researcher

\tableofcontents

% =============================================================
\mainmatter

% =============================================================
\part{Closure and Recovery: The Framework}

% =============================================================
\chapter{The Primitive Concept}

\epigraph{A map may be accurate in every detail it shows and still fail to
show the detail that matters most.}{}

\section{What This Book Is About}

Collapse preserves ontological embedding but not ontological recovery.

That sentence is the thesis of this book. Everything that follows is an
attempt to make it precise, to test it against the historical record, and to
ask what it implies about mathematical knowledge more broadly.

The sentence requires three concepts. \emph{Ontological embedding}: the old
objects survive as a recognizable subclass of a larger framework.
\emph{Ontological recovery}: the governing dynamics of the larger framework
remain confined to that subclass. \emph{Collapse}: the event by which the
invariant system of the old framework becomes globally unsatisfiable, forcing
enlargement.

None of these is standard terminology. They are introduced as names for a
structural phenomenon that keeps appearing across mathematics under different
local descriptions. The phenomenon is this. A theory develops an ontology,
a preferred class of objects and a set of governing operations, and that
ontology works reliably within a certain regime. Then the governing operations,
applied without restriction, produce objects the ontology cannot contain.
The ontology must be enlarged. But enlargement is not symmetric. The old
objects embed into the new framework as a recognizable subclass. The new
framework does not, in general, remain confined to the old subclass under
its own dynamics. The old theory survives. It does not remain sufficient.

\begin{remark}[On the word ``collapse'']
Throughout this book, \emph{collapse} is a technical term. Collapse does not
mean falsification, inconsistency, or catastrophic failure. It means precisely
the loss of recovery while embedding is retained. The old theory does not
become false. Its theorems remain valid within the embedded subclass. What
fails is the assumption that the original object class is sufficient for all
relevant purposes. Collapse is the formal event at which the old ontology
becomes a regime rather than a universe.
\end{remark}

\section{Continuation as Primitive}

Before introducing the formal apparatus, it is worth naming the conceptual
reorientation that the framework requires.

Traditional foundations of mathematics begin with objects: sets, points,
functions, distributions. Objects are the primitive substrate; relations and
operations are defined over them. The question is always: what are the things,
and what can we say about them?

The framework developed in this book begins differently. The primitive question
is not what a system is but whether an admissible continuation of it exists.
Given a state of a system, the primary datum is whether the governing operations
can be applied to produce a valid successor state within the current ontology.

Collapse is the event at which admissible continuation disappears.
Enlargement is the event at which admissible continuation is restored through
a larger ontology.

This reorientation has a formal expression in the admissible successor space
$\Admiss_B$ introduced below. It has a philosophical expression in the preface's
observation that objects appear fundamental only because they survive the
dynamics used to generate them. An ontology is not a collection of primitive
things. It is a collection of states that admit continuation under the
governing operations.

\section{Closure Operators and Dynamic Closure}

\begin{definition}[Dynamical Ontology]
A \emph{dynamical ontology} is a pair $(\Oold, \F)$ where $\Oold$ is a
class of objects and $\F = \{F_\alpha\}_{\alpha \in A}$ is a family of
governing operations, each $F_\alpha : \Oold \rightharpoonup \Oold$
(possibly partial).
\end{definition}

\begin{definition}[Dynamic Closure]
The \emph{closure} of a class $S \subseteq \Oold$ under $\F$ is
\[
  \cl_\F(S) = \bigcup_{n=0}^{\infty} \F^{(n)}(S),
\]
where $\F^{(0)}(S) = S$ and
$\F^{(n+1)}(S) = \bigcup_\alpha F_\alpha(\F^{(n)}(S))$. The class $S$ is
\emph{dynamically closed} under $\F$ if $\cl_\F(S) = S$.
\end{definition}

\begin{proposition}
$S$ is dynamically closed under $\F$ if and only if $F_\alpha(S) \subseteq S$
for all $\alpha \in A$.
\end{proposition}

\begin{proof}
If $F_\alpha(S) \subseteq S$ for all $\alpha$, then by induction
$\F^{(n)}(S) \subseteq S$ for all $n$, so $\cl_\F(S) \subseteq S$.
Since $S = \F^{(0)}(S) \subseteq \cl_\F(S)$, equality holds. Conversely,
if $\cl_\F(S) = S$ then $\F^{(1)}(S) \subseteq S$.
\end{proof}

\begin{example}[Smooth surfaces under area minimization]
Let $\Oold$ be smooth compact hypersurfaces with fixed boundary. In
dimensions $n \leq 7$, area-minimizing competitors of smooth surfaces are
smooth: $\Oold$ is dynamically closed. In dimension $n = 8$, the Simons
cone is an area-minimizing object that is not smooth: the operation escapes
$\Oold$. Dynamic closure fails.
\end{example}

\begin{example}[Real numbers under complex dynamics]
\label{ex:complex}
Let $\Oold = \Real \subset \mathbb{C} = \Onew$ and $F(z) = z^2 + i$.
For any $x \in \Real$, $F(x) = x^2 + i \notin \Real$. Embedding holds;
dynamic closure fails immediately.
\end{example}

\section{Embedding and Recovery}

\begin{definition}[Ontological Embedding]
\label{def:embedding}
A \emph{ontological embedding} is a structure-preserving injection
$\iota : \Oold \hookrightarrow \Onew$ between dynamical ontologies
$(\Oold, \F)$ and $(\Onew, \F')$ with $\F \subseteq \F'$.
\end{definition}

\begin{definition}[Ontological Recovery]
\label{def:recovery}
Given $\iota : \Oold \hookrightarrow \Onew$, the original ontology is
\emph{recoverable} within $(\Onew, \F')$ if
\[
  \F'(\Onew) \subseteq \iota(\Oold).
\]
\end{definition}

\begin{lemma}
Recovery implies that $\iota(\Oold)$ is dynamically closed under $\F'$.
\end{lemma}

\begin{proof}
If $\F'(\Onew) \subseteq \iota(\Oold)$, then in particular
$\F'(\iota(\Oold)) \subseteq \F'(\Onew) \subseteq \iota(\Oold)$.
\end{proof}

\begin{proposition}[Embedding Does Not Imply Recovery]
\label{prop:embedding-not-recovery}
There exist dynamical ontologies and embeddings for which $\iota(\Oold)$
is not dynamically closed under $\F'$.
\end{proposition}

\begin{proof}
Example~\ref{ex:complex}: $\Real \hookrightarrow \mathbb{C}$ under
$F(z) = z^2 + i$. For $x \in \Real$, $F(x) = x^2 + i \notin \Real$.
\end{proof}

\section{Strong and Weak Survival}

\begin{definition}[Weak and Strong Survival]
After enlargement $\iota : \Oold \hookrightarrow \Onew$:
\begin{itemize}
\item \emph{Weak survival} holds if $\Oold \hookrightarrow \Onew$.
\item \emph{Strong survival} holds if weak survival holds and every theorem
      $T$ valid universally in $\Oold$ remains valid universally in $\Onew$:
      $\Oold \models T \Rightarrow \Onew \models T$.
\end{itemize}
\end{definition}

\begin{theorem}
Strong survival implies weak survival. The converse fails.
\end{theorem}

\begin{proof}
Forward: immediate. Converse: smooth area-minimizing hypersurfaces embed in
rectifiable currents (weak survival), but the Simons cone is a current that
is not smooth (strong survival fails for the theorem ``area-minimizing objects
are smooth'').
\end{proof}

\begin{definition}[Regime Theorem]
A theorem $T$ that holds universally in $\Oold$ but only under additional
conditions in $\Onew$ is a \emph{regime theorem}. The conditions under which
it holds form the \emph{recoverable regime} for $T$.
\end{definition}

\begin{interlude}{Ontologies as Residues}
The objects of a theory appear fundamental only because they survive the
dynamics used to generate them. Smooth surfaces appear to be the natural
objects of variational geometry because they survive area minimization in
low dimensions. Classical probability distributions appear to be the natural
objects of inference because they survive Bayesian conditioning in commutative
settings. Semantic scopes appear to be the natural objects of meaning because
they survive admissible resolution in well-structured discourse.

In each case, the apparent fundamentality of the object class is a feature
of the recoverable regime, not of the world. The objects are what persist.
They are residues of the dynamics, not its foundations.

This observation inverts the usual direction of mathematical ontology.
Traditionally, one begins with objects and derives dynamics. The framework
developed here suggests the opposite order: dynamics come first, and objects
are what they leave invariant.
\end{interlude}

\section{Invariant Systems and Collapse}

\begin{definition}[Invariant System]
An \emph{invariant system} $\mathcal{V} = \{V_1, \ldots, V_k\}$ for
$(\Oold, \F)$ is a collection of properties jointly required of all
admissible objects.
\end{definition}

\begin{definition}[Admissible Successor Space and Inadmissibility Pressure]
\[
  \Admiss_B = \{ B' \in \Oold \mid B' \text{ satisfies all }
  V_i \in \mathcal{V} \}.
\]
\[
  \Inadmiss(B, F) = \dist(F(B),\, \Admiss_B).
\]
The admissible successor space is defined by the invariant system, not by
which operations succeed. This prevents circularity.
\end{definition}

\begin{definition}[Refuse and Collapse]
\label{def:refuse-collapse}
An operation $F$ is \emph{refused} at $B$ if $\Inadmiss(B,F) \geq \theta > 0$
while $\Admiss_B \neq \varnothing$: local failure, global ontology intact.

The ontology \emph{collapses} when $\Admiss_B = \varnothing$: no object in
$\Oold$ satisfies the full invariant system. Admissible continuation has
ceased to exist within the current ontology.
\end{definition}

\section{Obstruction Theory}

\begin{definition}[Obstruction Set]
\[
  \W(\Oold, \Onew, \F') =
  \{ w \in \Onew \setminus \iota(\Oold) \mid w \in \F'(\Onew) \}.
\]
An element $w \in \W$ is an \emph{obstruction}: an object produced by the
governing operations that belongs to the enlarged ontology but not to the
original one.
\end{definition}

\begin{theorem}[Obstruction Implies Non-Recovery]
\label{thm:obstruction-nonrecovery}
If $\W \neq \varnothing$ then recovery fails: $\F'(\Onew) \not\subseteq
\iota(\Oold)$.
\end{theorem}

\begin{proof}
Take $w \in \W$. Then $w \in \F'(\Onew)$ and $w \notin \iota(\Oold)$.
\end{proof}

\begin{remark}[Obstruction versus Collapse]
Theorem~\ref{thm:obstruction-nonrecovery} establishes non-recovery from an
obstruction. Non-recovery and collapse are related but distinct. Collapse
($\Admiss_B = \varnothing$) concerns the original ontology's inability to
provide admissible continuations. Non-recovery concerns the enlarged ontology's
inability to remain confined to the original. In all three main case studies,
the invariant system includes dynamic closure as a requirement, so the two
conditions coincide. We will make this explicit in each case.
\end{remark}

\begin{theorem}[Irreversibility of Enlargement]
\label{thm:irreversibility}
Suppose $\iota : \Oold \hookrightarrow \Onew$, recovery fails, and
$\cl_{\F'}(\iota(\Oold)) = \Onew$ (the enlarged ontology is reachable from
the embedded subclass). Then $\iota(\Oold)$ is not dynamically closed under
$\F'$.
\end{theorem}

\begin{proof}
Since recovery fails, $\F'(\Onew) \not\subseteq \iota(\Oold)$. By
reachability, $\Onew = \cl_{\F'}(\iota(\Oold))$, so there exists a finite
chain $x_0 \in \iota(\Oold)$, $x_1 = F'_{\alpha_1}(x_0)$, \ldots,
$x_n = F'_{\alpha_n}(x_{n-1})$ with $x_n \notin \iota(\Oold)$ for the
first such $n$. Then $F'_{\alpha_n}(x_{n-1}) \notin \iota(\Oold)$ while
$x_{n-1} \in \iota(\Oold)$: $\iota(\Oold)$ is not closed under
$F'_{\alpha_n}$.
\end{proof}

\begin{definition}[Obstruction Accessibility]
An obstruction $w \in \W$ is \emph{accessible} if there exists a feasible
procedure $P$ such that $P$ produces $w$ or a witness to $w$'s existence.
Accessibility is relative to available tools and changes as methods develop.
\end{definition}

\section{Closure as Reflection: A Categorical View}

The framework has a natural expression in category theory.

\begin{definition}[Reflective Subcategory]
A subcategory $\bO \subset \bOn$ is \emph{reflective} if the inclusion
functor $\iota : \bO \hookrightarrow \bOn$ has a left adjoint
$R : \bOn \to \bO$, called the \emph{reflector}. The adjunction
$R \dashv \iota$ gives natural isomorphisms
$\mathrm{Hom}_{\bO}(R(X), A) \cong \mathrm{Hom}_{\bOn}(X, \iota(A))$
for all $X \in \bOn$, $A \in \bO$.
\end{definition}

Reflective subcategories are precisely those that are ``closed under the
operations of the larger category in a recoverable way.'' The reflector
provides a canonical projection back into the small category.

\begin{proposition}[Recovery as Reflection]
Ontological recovery holds if and only if the inclusion
$\iota : \bO \hookrightarrow \bOn$ participates in a reflective adjunction
$R \dashv \iota$ such that $R$ commutes with the governing operations:
$R \circ F' \cong F \circ R$.
\end{proposition}

\begin{corollary}[Collapse as Failure of Reflection]
\label{cor:collapse-reflection}
Ontological collapse occurs when no reflector $R : \bOn \to \bO$ compatible
with the governing operations exists. There is no canonical way to project
outputs of $\F'$ back into $\bO$.
\end{corollary}

\begin{example}[GMT in Categorical Terms]
Let $\bO$ be smooth minimal hypersurfaces and $\bOn$ integer-multiplicity
rectifiable currents. The inclusion does not have a reflector compatible with
area minimization in dimension eight: there is no canonical way to smooth out
the Simons cone while preserving its minimizing property. Reflection fails;
collapse is confirmed.
\end{example}

\section{The Six-Stage Pattern}

\begin{enumerate}[label=\textbf{Stage \arabic*.}, leftmargin=*, itemsep=5pt]

\item \textbf{Operational Failure (Refuse).}
$\Inadmiss(B, F) \geq \theta$ while $\Admiss_B \neq \varnothing$.
Operation blocked locally; ontology intact globally.

\item \textbf{Ontological Collapse.}
$\Admiss_B = \varnothing$. Admissible continuation has ceased to exist
within the current ontology. This is an objective mathematical fact,
independent of recognition.

\item \textbf{Enlargement Availability.}
$\Onew \supset \iota(\Oold)$ is constructed with $\F'(\Onew) \subseteq \Onew$.
May precede recognized collapse.

\item \textbf{Obstruction Existence.}
$\W \neq \varnothing$. A specific object certifies that recovery fails.

\item \textbf{Obstruction Accessibility.}
A feasible procedure produces a witness. The gap between Stages~4 and~5
is the \emph{accessibility gap}, governed by the structure of the obstruction
region relative to available tools --- not by the sociology of the field.

\item \textbf{Ontological Commitment.}
The field accepts the enlarged ontology as necessary. Universal claims of
the old theory are reinterpreted as regime theorems.

\end{enumerate}

\begin{metatheorem}[Closure-Recovery Principle]
\label{meta:CRP}
Let $(\Oold, \F)$ be a dynamical ontology. If $\cl_\F(\Oold) \neq \Oold$,
then either $\F$ must be restricted or $\Oold$ must be enlarged.
Historically, enlargement is the repeated outcome.
\end{metatheorem}

% =============================================================
\part{The Geometry Case: A Completed Collapse}

% =============================================================
\chapter{Functionals and the First Enlargement}

\epigraph{Bernoulli asked for the fastest curve.\\
Lagrange realized the question was about all curves at once.}{}

\section{The Brachistochrone and the Birth of Functionals}

In 1696 Johann Bernoulli posed a problem that looks, at first, like an
exercise in mechanics. A bead slides without friction along a wire connecting
two points in a vertical plane. What shape should the wire take so that the
bead reaches the lower point in the shortest possible time?

The answer is a cycloid. But Bernoulli understood that the significant
achievement was not the answer but the method. To find the fastest wire,
one does not examine a particular wire. One examines all wires simultaneously
and asks which is stationary with respect to travel time. Travel time is not
a function of a real variable; it is a function of an entire curve: a
\emph{functional}. The problem of finding curves making a functional stationary
is the subject of the calculus of variations.

Let the wire be described by $y(x)$ with $y(x_0) = y_0$, $y(x_1) = y_1$.
Conservation of energy gives speed $v = \sqrt{2g(y_0 - y)}$ at height $y$.
The travel time is
\[
  T[y] = \int_{x_0}^{x_1}
  \frac{\sqrt{1 + (y')^2}}{\sqrt{2g(y_0 - y)}}\, dx.
\]
This is a functional: a map from the space of curves to the real numbers.

\section{The Euler-Lagrange Equation}

\begin{theorem}[Euler-Lagrange]
\label{thm:EL}
If $y$ is $C^2$, satisfies fixed boundary conditions $y(a) = A$, $y(b) = B$,
and $\delta J[y] = 0$ for all compactly supported smooth perturbations $\eta$
vanishing at the endpoints, where $J[y] = \int_a^b L(x, y, y')\, dx$, then
\[
  \frac{\partial L}{\partial y}
  - \frac{d}{dx}\frac{\partial L}{\partial y'} = 0.
\]
\end{theorem}

\begin{proof}
Set $y_\varepsilon = y + \varepsilon\eta$. Differentiating at $\varepsilon = 0$:
\[
  \frac{d}{d\varepsilon}\bigg|_0 J[y_\varepsilon]
  = \int_a^b \Bigl(\frac{\partial L}{\partial y}\eta
    + \frac{\partial L}{\partial y'}\eta'\Bigr) dx = 0.
\]
Integrating the second term by parts and using $\eta(a) = \eta(b) = 0$:
\[
  \int_a^b \Bigl(\frac{\partial L}{\partial y}
  - \frac{d}{dx}\frac{\partial L}{\partial y'}\Bigr)\eta\, dx = 0.
\]
Since $\eta$ is arbitrary, the integrand vanishes everywhere.
\end{proof}

\begin{remark}[The first ontological enlargement]
The passage from points to functionals is itself an ontological enlargement.
The brachistochrone could not be stated without the concept of a functional.
Functions on curves embed in the space of all functions on an interval,
but the variational calculus operates on function spaces whose governing
dynamics do not remain confined to finite-dimensional subspaces. Embedding
holds; recovery fails immediately. The calculus of variations is the enlarged
ontology that restored closure, at the cost of requiring the Euler-Lagrange
equation rather than simple differentiation.
\end{remark}

\section{The Plateau Problem and Minimal Surfaces}

For a surface expressed as a graph $z = f(x,y)$, the minimal surface equation
(vanishing mean curvature $H = 0$) is
\[
  (1 + f_y^2)f_{xx} - 2f_x f_y f_{xy} + (1 + f_x^2)f_{yy} = 0.
\]
Physical soap films satisfy this equation: area minimization produces
vanishing mean curvature. The Plateau problem asks whether, given any closed
curve, a minimal surface spanning it exists.

\section{Bernstein's Theorem and the Illusion of Universality}

\begin{theorem}[Bernstein, 1915 \cite{bernstein1915}]
Let $f : \Real^2 \to \Real$ be $C^2$ with minimal graph. Then $f$ is affine:
the only complete minimal graphs in $\Real^3$ are planes.
\end{theorem}

This was extended through dimension seven by De~Giorgi \cite{degiorgi1965},
Almgren \cite{almgren1966}, and Simons \cite{simons1968}:

\begin{theorem}[Bernstein in Low Dimensions]
The Bernstein theorem holds for entire minimal hypergraphs $f:\Real^{n-1}
\to \Real$ for $n \leq 8$: any $C^2$ solution of the minimal surface equation
defined on all of $\Real^{n-1}$ must be linear.
\end{theorem}

With each dimensional extension the expectation hardened: minimization
enforces simplicity universally. The invariant system (smooth, minimizing,
globally defined) appeared jointly satisfiable in all dimensions. The collapse
was present in dimension eight long before it was recognized.

% =============================================================
\chapter{The Simons Cone and Obstruction}

\epigraph{The cone is not important because it is singular.\\
It is important because it is necessary.}{}

\section{The Construction}

In $\Real^8$ define
\[
  \mathcal{C} = \{ (x_1,\ldots,x_8) \in \Real^8 \mid
  x_1^2 + x_2^2 + x_3^2 + x_4^2 = x_5^2 + x_6^2 + x_7^2 + x_8^2 \}.
\]
This is the \emph{Simons cone} \cite{simons1968}: a seven-dimensional cone
with vertex at the origin, smooth everywhere except the origin.

\begin{theorem}[Simons 1968; Bombieri-De~Giorgi-Giusti 1969
  \cite{simons1968,bombieri1969}]
\label{thm:simons}
The Simons cone $\mathcal{C} \subset \Real^8$ is area-minimizing. Moreover,
the Bernstein theorem fails for $n \geq 9$: there exist complete minimal
hypergraphs in $\Real^n$ that are not affine hyperplanes.
\end{theorem}

Dimension eight is the sharp boundary.

\section{Dimension as Control Parameter}

The mechanism is competition between curvature cost and area gain. In low
dimensions, singularities are energetically penalized: curvature concentrates
at a singular point at greater cost than the area saved. Additional dimensions
distribute curvature across more principal directions, weakening the penalty.
At dimension eight, the balance tips.

Stability is established by the second variation:

\begin{definition}[Second Variation and Stability]
For smooth minimal $\Sigma$ and compactly supported $\phi$:
\[
  \delta^2\mathcal{A}(\Sigma)[\phi,\phi]
  = \int_\Sigma (|\nabla\phi|^2 - |A|^2\phi^2)\, d\mathcal{H}^{n-1}.
\]
$\Sigma$ is \emph{stable} if $\delta^2\mathcal{A}[\phi,\phi] \geq 0$ for
all compactly supported $\phi$.
\end{definition}

Simons showed that $\mathcal{C}$ satisfies the estimates needed for stability
in $\Real^8$ but not in lower-dimensional analogues. The balance between
$|\nabla\phi|^2$ and $|A|^2\phi^2$ tips at precisely the ambient dimension
where the cone lives.

\section{The Cone as Obstruction}

In the language of Chapter~1, identifying each component:

\emph{Original ontology $\Oold$}: smooth minimal hypersurfaces.
\emph{Governing operation}: area minimization.
\emph{Invariant system}: $V_1$ smoothness, $V_2$ area-minimization,
$V_3$ global existence for all boundary data.
\emph{Enlarged ontology $\Onew$}: integer-multiplicity rectifiable currents.
\emph{Obstruction}: the Simons cone $\mathcal{C} \in \Onew \setminus
\iota(\Oold)$, produced by area minimization, certifying that recovery fails.

By Theorem~\ref{thm:obstruction-nonrecovery}, the cone certifies non-recovery.
Since the invariant system includes dynamic closure, non-recovery and collapse
coincide: the smooth-surface ontology has no admissible continuation for the
cone's boundary conditions.

The cone does not create the failure. It certifies a collapse already present
in the mathematics of $\Real^8$. The cone is what makes the collapse
impossible to ignore.

\section{Smoothness Becomes a Regime}

After Theorem~\ref{thm:simons}, ``area-minimizing surfaces are smooth'' is
reinterpreted as a regime theorem. It holds in dimensions $\leq 7$ and as
a statement about the regular part of minimizing currents in higher dimensions.

\begin{theorem}[Regularity, Sharp Form \cite{federer1969,simon1983}]
For an area-minimizing integer-multiplicity rectifiable $n$-current in
$\Real^{n+1}$:
\begin{enumerate}[label=(\arabic*)]
\item The regular set $\mathrm{Reg}(T)$ is open and dense in $\supp(T)$.
\item $\dim_{\mathcal{H}}\mathrm{Sing}(T) \leq n - 7$.
\item In dimensions $n \leq 6$: $\mathrm{Sing}(T) = \varnothing$.
\item In dimension $n = 7$: $\mathrm{Sing}(T)$ consists of isolated points.
\end{enumerate}
\end{theorem}

The number seven appears because it is the lowest dimension in which the
Simons cone (a dimension-seven object) can be a stable minimizer. Below the
threshold, singular minimizers cannot persist; above it, they can.

% =============================================================
\chapter{Geometric Measure Theory as Enlarged Ontology}

\epigraph{Classical differential geometry assumes smooth surfaces.\\
Geometric measure theory asks what minimization actually produces.}{}

\section{The Closure Failure of Smooth Surfaces}

Minimizing sequences of smooth surfaces need not have convergent subsequences
in the smooth category. Curvatures can concentrate; singularities can develop;
topological type can change. The class of smooth surfaces is not closed under
sequential limits, which variational methods require.

\section{Currents: Closure Restored}

\begin{definition}[Current and Mass \cite{federer1969}]
A \emph{$k$-dimensional current} in $\Real^n$ is a continuous linear
functional $T : \Omega^k_c(\Real^n) \to \Real$, with boundary
$\partial T(\omega) = T(d\omega)$ and mass
$\mathbf{M}(T) = \sup\{T(\omega) : \sup|\omega| \leq 1\}$.

An \emph{integer-multiplicity rectifiable $k$-current} is
$T(\omega) = \int_M \langle\omega, \xi\rangle\theta\, d\mathcal{H}^k$
where $M$ is $k$-rectifiable, $\xi$ is an orientation, and
$\theta : M \to \mathbb{Z}_{>0}$.
\end{definition}

Smooth oriented surfaces are the special case $\theta \equiv 1$.
The Simons cone is a rectifiable current with a singularity at the origin.
Both live in the same enlarged ontology.

\begin{theorem}[Federer-Fleming Compactness \cite{federer1969}]
\label{thm:FF}
A sequence of integer-multiplicity rectifiable $k$-currents with uniformly
bounded mass and boundary mass has a weakly convergent subsequence to an
integer-multiplicity rectifiable $k$-current.
\end{theorem}

\begin{remark}[Closure restored]
Theorem~\ref{thm:FF} is precisely the statement that the enlarged ontology
is dynamically closed under the sequential limits that variational methods
require. Minimizing sequences that escape the smooth category are contained
within the current category. The governing operation closes on the enlarged
ontology.
\end{remark}

\begin{theorem}[Existence of Area-Minimizing Currents]
For any rectifiable $(k-1)$-cycle $\Gamma$, there exists an integer-multiplicity
rectifiable $k$-current $T$ with $\partial T = \Gamma$ achieving the minimum
mass.
\end{theorem}

\section{Tangent Cones and Local Models}

\begin{definition}[Tangent Cone]
Let $T$ be a current, $x \in \mathrm{Sing}(T)$. The \emph{tangent cone} at $x$
is the weak limit as $r \to 0$ of the rescaled currents
$T_{x,r}(\omega) = r^{-n}T((\tau_x \circ \delta_{1/r})_\#\omega)$.
\end{definition}

Tangent cones are local models for singularities. The Simons cone appears
as the tangent cone for a generic class of singularities in dimension eight.
It is an attractor in the space of rescaled currents: a renormalization fixed
point. This explains why it is not merely a counterexample but a structural
object. Singularities, in the vicinity of this particular renormalization
fixed point, look like cones.

\begin{interlude}{Return and Recovery}
The purpose of enlargement is not escape from the original ontology.
It is the recovery of those conditions under which the original ontology
remains valid.

Regularity theory is a theory of return. Given a minimizing current in the
enlarged ontology, regularity theory determines where and to what extent the
original smooth description is recoverable. The regular set $\mathrm{Reg}(T)$
is the region of return: where the object behaves as the old ontology
predicted. The singular set $\mathrm{Sing}(T)$ is where return is impossible.

This structure repeats across all the case studies in this book. Asymptotic
frequentism in quantum inference is a theorem of return: the posterior
eventually reaches the true state, even if the path is non-monotone.
Scope completion in Spherepop is a theorem of return: the load eventually
reaches zero, even if intermediate operations temporarily increase it.

Return to the old ontology is possible in the generic case. The enlarged
ontology does not eliminate the old one; it situates it. It explains when
the old description holds and when it does not. The enlarged framework is
valuable precisely because it provides a theory of return.
\end{interlude}

\section{The GMT Case: Five Claims}

The GMT case establishes five claims that the subsequent case studies
reproduce in modified form.

Collapse is objective and precedes recognition. The invariant system was
unsatisfiable in dimension eight before the Simons cone was constructed.

The enlarged ontology must be available before an obstruction can be
recognized. GMT existed before the cone.

The accessibility gap is structural, not sociological. The cone was
constructible with available tools within approximately a decade.

Strong survival fails while weak survival holds. The universal smoothness
theorem becomes a regime theorem.

The enlarged ontology has its own structure, richer than the original.
Regularity theory, tangent cones, mean curvature flow: these research programs
do not exist within the smooth category. The enlargement opens new territory.

% =============================================================
\part{The Inference Case: An Ongoing Collapse}

% =============================================================
\chapter{Classical Bayesian Inference and Its Geometry}

\epigraph{Probability is not a property of the world. It is a property of
our relationship to the world. But that relationship has a geometry.}{}

\section{Inference as Dynamical Process}

The original ontology is the class of probability distributions $\Oold =
\{\pi : \Theta \to \Real_{\geq 0} \mid \int\pi = 1\}$ over a hypothesis
space $\Theta$. The governing operation is Bayesian conditioning:
$F_e(\pi)(h) = \pi(h \mid e) = \pi(e \mid h)\pi(h)/\pi(e)$. The invariant
system consists of $V_1$ (classical joint state existence), $V_2$ (unique
conditioning), and $V_3$ (monotone improvement under repetition). All three
are jointly satisfiable in the commutative setting.

\section{Kullback-Leibler Divergence and the Fisher Metric}

\begin{definition}[Kullback-Leibler Divergence]
$\KL(p \| q) = \int p(x)\log(p(x)/q(x))\, dx \geq 0$, with equality iff
$p = q$ almost everywhere.
\end{definition}

\begin{definition}[Fisher Information Metric]
$g_{ij}(\theta) = \mathbb{E}_{p_\theta}
[(\partial_i\log p_\theta)(\partial_j\log p_\theta)]$.
\end{definition}

The Fisher metric is, up to rescaling, the unique Riemannian metric on
statistical manifolds invariant under sufficient statistics
\cite{chentsov1982}. Bayesian conditioning is an information projection in
this geometry: the posterior minimizes relative entropy to the prior subject
to matching the expected sufficient statistics of the evidence.

\section{The Data Processing Inequality and No Epistemic Harm}

\begin{theorem}[Data Processing Inequality]
For any measurable $\phi : \mathcal{X} \to \mathcal{Y}$ and distributions $p$, $q$:
$\KL(\phi_* p \| \phi_* q) \leq \KL(p \| q)$.
\end{theorem}

\begin{proof}[Sketch]
Chain rule: $\KL(p_{XY} \| q_{XY}) = \KL(p_X \| q_X) +
\KL(p_{Y|X} \| q_{Y|X})$. Processing by $\phi$ marginalizes the
non-negative second term.
\end{proof}

\begin{corollary}[No Epistemic Harm, Classical]
\label{cor:NEH}
Let $\pi^*$ be the true distribution, $\pi_n$ the posterior after $n$
observations. Then $\KL(\pi^* \| \pi_{n+1}) \leq \KL(\pi^* \| \pi_n)$.
Repeated Bayesian updating weakly decreases divergence from truth.
\end{corollary}

This is invariant $V_3$: monotone improvement under the governing operation.

\section{Asymptotic Frequentism}

\begin{theorem}[Bernstein-von Mises]
Under mild regularity conditions, the posterior $\pi_n(\cdot \mid x_1,\ldots,
x_n)$ converges in total variation to $\delta_{\theta^*}$ as $n \to \infty$,
at a rate asymptotically independent of the prior $\pi_0$.
\end{theorem}

The true distribution is the attractor; the prior is the initial condition.
In the classical case, no epistemic harm (local monotonicity) and asymptotic
frequentism (global convergence) both hold. In the quantum case, they may
come apart.

\section{Schr\"{o}dinger Bridges and Process Updating}

\begin{definition}[Schrödinger Bridge \cite{schrodinger1932}]
Let $\mathbb{P}$ be Wiener measure and $\mu_0$, $\mu_1$ measures on $\Real^d$.
The \emph{Schrödinger bridge} is
\[
  \mathbb{Q}^* = \operatorname{argmin}\{
  \KL(\mathbb{Q}\|\mathbb{P}) \mid \mathbb{Q}_0 = \mu_0,\,
  \mathbb{Q}_1 = \mu_1\}.
\]
\end{definition}

The Sinkhorn algorithm \cite{sinkhorn1967} computes Schrödinger bridges by
alternately projecting onto the marginal constraints. This is inference at
the process level: instead of updating a belief state, one updates a coupling.
The enlarged ontology of processes embeds the original ontology of states;
whether recovery holds is the question the quantum case approaches.

% =============================================================
\chapter{Quantum Inference and the Prospective Collapse}

\epigraph{Quantum mechanics is not probability theory with complex numbers.\\
It is probability theory with a different geometry.}{}

\section{Density Matrices and Noncommutativity}

\begin{definition}[Density Matrix and Quantum Relative Entropy]
A \emph{density matrix} on $\Hilb$ is $\rho \in \mathcal{L}(\Hilb)$ with
$\rho \geq 0$ and $\tr(\rho) = 1$. For density matrices with
$\supp(\rho) \subseteq \supp(\sigma)$:
\[
  D(\rho\|\sigma) = \tr(\rho\log\rho - \rho\log\sigma) \geq 0.
\]
\end{definition}

The fundamental structural difference: density matrices generically fail
to commute, $[\rho,\sigma] \neq 0$. This destroys the canonical joint-state
structure classical inference requires.

\begin{theorem}[Quantum Data Processing Inequality \cite{wilde2015}]
For any CPTP map $\mathcal{E}$: $D(\mathcal{E}(\rho)\|\mathcal{E}(\sigma))
\leq D(\rho\|\sigma)$.
\end{theorem}

\section{The Petz Recovery Map}

\begin{definition}[Petz Recovery Map \cite{petz1988}]
Let $\sigma \in \density(\Hilb_{AB})$ with marginal $\sigma_B = \tr_A(\sigma)$.
The \emph{Petz recovery map} is
\[
  \mathcal{R}_{\sigma,\tr_A}(X)
  = \sigma^{1/2}(\sigma_B^{-1/2} X \sigma_B^{-1/2} \otimes \mathbf{1}_A)
    \sigma^{1/2}.
\]
\end{definition}

\begin{proposition}[Properties of the Petz Map]
\begin{enumerate}[label=(\arabic*)]
\item $\mathcal{R}_{\sigma,\tr_A}$ is CPTP.
\item Quantum detailed balance:
      $D(\rho\|\sigma) = D(\tr_A(\rho)\|\sigma_B) +
       D(\rho\|\mathcal{R}_{\sigma,\tr_A}(\tr_A(\rho)))$.
\item In the commutative regime, it reduces to classical Bayesian conditioning.
\end{enumerate}
\end{proposition}

Property~(3) confirms $\Oold \hookrightarrow \Onew$: classical inference
embeds in quantum inference. Property~(2) is the quantum Bayes rule.
The question is recovery.

\section{A Worked Petz Recovery Example}

Let
\[
  \rho = \begin{pmatrix} 0.7 & 0.25 \\ 0.25 & 0.3 \end{pmatrix},
  \qquad
  \sigma = \begin{pmatrix} 0.6 & 0 \\ 0 & 0.4 \end{pmatrix}.
\]
These matrices do not commute: $[\rho,\sigma] \neq 0$, placing the system
outside the classical regime. Consider the dephasing channel
$\mathcal{E}(X) = \mathrm{diag}(X_{11}, X_{22})$:
\[
  \mathcal{E}(\rho) = \begin{pmatrix}0.7 & 0 \\ 0 & 0.3\end{pmatrix},
  \qquad \mathcal{E}(\sigma) = \sigma.
\]
The quantum data-processing inequality guarantees
$D(\mathcal{E}(\rho)\|\mathcal{E}(\sigma)) \leq D(\rho\|\sigma)$.
Information is lost under coarse-graining, as expected.

The Petz recovery map produces a recovered state
$\rho_1 = \mathcal{R}_{\sigma,\mathcal{E}}(\mathcal{E}(\rho))$.
The diagnostic quantity is
\[
  \Delta_n = D(\rho^*\|\rho_{n+1}) - D(\rho^*\|\rho_n),
\]
where $\rho^*$ is the target state and $\rho_n$ the state after $n$ Petz
recovery cycles.

Classical Bayesian inference predicts $\Delta_n \leq 0$ for every $n$
(Corollary~\ref{cor:NEH}). The conjectural obstruction regime predicts the
existence of specific triples $(\rho, \sigma, \mathcal{E})$ and targets
$\rho^*$ for which $\Delta_n > 0$ for some finite $n$, even though
$D(\rho^*\|\rho_n) \to 0$ as $n \to \infty$.

Such a trajectory temporarily moves away from the target before ultimately
converging. This is not a failure of inference. It is a failure of monotonicity:
the failure of $V_3$ while global convergence (the attractor structure) may
survive. The obstruction is not non-convergence. The obstruction is the
failure of monotonicity itself, which would certify that the classical
invariant system is unsatisfiable in the noncommutative regime.

The GMT parallel is exact. Minimizing currents exist and generically converge
to area-minimizing limits (global existence survives); but not every
minimizing current is smooth (local regularity fails). Quantum inference may
generically converge to the true state (attractor structure survives); but
not every inference trajectory is monotone (local monotonicity fails).

\section{The Invariant System and What Collapses}

The classical invariant system for quantum inference:

$V_1$: a classical joint state encodes the prior-evidence relationship.\\
$V_2$: conditioning is uniquely determined by the joint state.\\
$V_3$: $D(\rho^*\|\rho_{n+1}) \leq D(\rho^*\|\rho_n)$ for all $n$.

Under noncommutativity, $V_1$ and $V_2$ are structurally unsatisfiable:
quantum systems generically lack classical joint states. Whether $V_3$
fails is the open question that determines whether the prospective collapse
is established.

\begin{conjecture}[Noncommutativity Threshold]
\label{conj:threshold}
There exists a measure $\eta$ of effective noncommutativity and a critical
value $\eta_c$ such that: for $\eta < \eta_c$, the Petz map satisfies $V_3$;
for $\eta > \eta_c$, violations of $V_3$ are generic; the transition is sharp.
\end{conjecture}

This is a falsifiable prediction. Gradual degradation of monotonicity across
all $\eta$ would weaken the Simons cone analogy considerably. Sharp threshold
behavior would confirm the structural parallel quantitatively.

\section{Established versus Prospective Collapse}

\begin{definition}[Established and Prospective Collapse]
A collapse is \emph{established} if there is a formal proof that
$\Admiss_B = \varnothing$ in $\Oold$ for the relevant invariant system.
A collapse is \emph{prospective} if strong structural and numerical evidence
exists but formal proof is absent.
\end{definition}

The GMT case is established. The quantum Bayes case is prospective.
This distinction does not weaken the framework; it makes it diagnostic
rather than merely retrospective. A framework that only classifies completed
revolutions is a catalogue.

% =============================================================
\part{The Formal Case: Spherepop as Abstract Model}

% =============================================================
\chapter{Scope Graphs and Semantic Topology}

\epigraph{A bubble is not a thing. It is a commitment deferred.}{}

\section{The Spherepop Calculus: Definitions}

\begin{definition}[Semantic Bubble]
A \emph{semantic bubble} is a triple $B = (C, E, U)$ where $C$ is a
contextual binding set, $E$ is an expectation structure, and $U(B) \geq 0$
is the unresolved semantic load. $B$ is \emph{open} if $U(B) > 0$,
\emph{closed} if $U(B) = 0$.
\end{definition}

\begin{definition}[Containment Structure, Scope Stack, Load Functional]
A \emph{containment structure} is $(\mathcal{B}, \prec)$ with $\prec$ a
strict partial order. The \emph{scope stack} is $\Sigma = [B_1 \prec \cdots
\prec B_n]$. The \emph{semantic load} is
$L(\Sigma) = \sum_i w_i U(B_i)$ where $w_i = f(d_i, s_i, r_i, c_i)$
weights depth, salience, recency, and causal centrality.
\end{definition}

\section{Scope Graphs}

Containment trees are too restrictive for cross-cutting dependencies.
Human reasoning, mathematical proof, software architecture, and scientific
explanation frequently involve structures where a scope depends on another
without being spatially contained within it.

\begin{definition}[Scope Graph]
A \emph{scope graph} is a finite directed graph $G = (V, E)$ with semantic
load $U : V \to \Real_{\geq 0}$. Vertices are semantic bubbles; edge
$B_i \to B_j$ means admissibility of $B_i$ depends on resolution of $B_j$.
\end{definition}

Containment trees appear as scope graphs with at most one outgoing dependency
per vertex and no directed cycles. The scope graph separates logical dependence
from spatial containment.

\begin{definition}[Dependency Closure]
The \emph{dependency closure} $\cl(G)$ is the smallest graph containing $G$
and closed under transitivity: $(B_i, B_j) \in \cl(G)$ whenever there exists
a directed path from $B_i$ to $B_j$.
\end{definition}

\begin{proposition}
$\cl(\cl(G)) = \cl(G)$.
\end{proposition}

\begin{proof}
The first closure adjoins every transitive dependency. A second application
introduces no additional edges.
\end{proof}

Dependency closure is the Spherepop analogue of the topological closure
operator: it closes a semantic structure under dependency implication rather
than a set under limits.

\section{Collapse Cycles as Obstructions}

\begin{definition}[Collapse Cycle]
A \emph{collapse cycle} is a directed cycle $B_1 \to B_2 \to \cdots \to
B_n \to B_1$ in which every edge represents an unresolved dependency.
\end{definition}

\begin{theorem}[Collapse Cycle Theorem]
\label{thm:collapse-cycle}
If a scope graph contains a collapse cycle then $\Admiss = \varnothing$.
\end{theorem}

\begin{proof}
Resolution of $B_1$ requires prior resolution of $B_2$. Resolution of $B_2$
requires prior resolution of $B_3$. By induction, resolution of $B_n$
requires prior resolution of $B_1$. No node may be resolved first.
Therefore no admissible Pop sequence exists and $\Admiss = \varnothing$.
\end{proof}

The collapse cycle plays the role of the Simons cone in the graph formalism.
It does not create the collapse; it certifies it. A collapse cycle is a stable
dependency configuration that makes the collapse impossible to ignore,
exactly as the Simons cone is a stable singular configuration certifying that
smooth minimizers do not exist for the cone's boundary conditions.

\section{Minimal Obstructions}

\begin{definition}[Minimal Obstruction]
A subgraph $H \subseteq G$ is a \emph{minimal obstruction} if $H$ satisfies
Theorem~\ref{thm:collapse-cycle} (it contains a collapse cycle) but every
proper subgraph of $H$ admits an admissible resolution sequence.
\end{definition}

\begin{theorem}[Existence of Minimal Obstructions]
Every finite collapsed scope graph contains a minimal obstruction.
\end{theorem}

\begin{proof}
Let $\mathcal{O} = \{H \subseteq G : H \text{ is collapsed}\}$. This is
nonempty since $G \in \mathcal{O}$. Since $G$ is finite, $\mathcal{O}$
has an element of minimal cardinality. Any such element is a minimal
obstruction: removing any edge or vertex would restore admissibility.
\end{proof}

This establishes an obstruction theory for Spherepop analogous to
forbidden-minor theory in graph topology. Minimal obstructions are the
smallest witnesses to collapse: the Spherepop analogues of the Simons cone.

\section{Semantic Entropy and the Obstruction Measure}

\begin{definition}[Semantic Entropy]
The \emph{semantic entropy} of a scope graph $G$ is
\[
  S(G) = \log|\Admiss(G)|,
\]
the log-cardinality of admissible resolution sequences ($-\infty$ when
$\Admiss = \varnothing$, i.e., at collapse).
\end{definition}

High $S(G)$ means many admissible continuations: the system is far from
collapse. Low $S(G)$ means few: collapse is imminent. $S(G) = -\infty$
means collapse has occurred. Semantic entropy is the Spherepop analogue
of $\dist(F(B), \Admiss_B)$.

\begin{theorem}[Semantic Thermodynamics]
Under pure Pop sequences, semantic entropy is non-decreasing. Under Meld
and Reframe, semantic entropy may temporarily decrease before stabilizing.
\end{theorem}

\begin{proof}[Sketch]
Each admissible Pop closes one scope, removing one constraint from the
parent's admissibility condition, weakly increasing $|\Admiss|$. Meld and
Reframe introduce new compatibility constraints that may temporarily
reduce $|\Admiss|$ before they stabilize.
\end{proof}

\section{The Collapse Quotient and Irreversibility}

\begin{definition}[Collapse Quotient]
Let $\sim_\rho$ be the equivalence relation on bubble histories: two histories
are equivalent when they produce the same admissible semantic representative
under all available resolution sequences. The \emph{collapsed residue} of
$B$ is $B/{\sim_\rho}$.
\end{definition}

\begin{construction}[Universal Property]
The quotient map $q : B \to B/{\sim_\rho}$ satisfies: for any $f : B \to X$
constant on $\sim_\rho$-classes, there exists unique $\bar{f} : B/{\sim_\rho}
\to X$ with $f = \bar{f} \circ q$.
\end{construction}

\begin{proposition}[Irreversibility of Spherepop Collapse]
\label{prop:spherepop-irrev}
The quotient map $q : B \to B/{\sim_\rho}$ does not generally admit a section.
Generative history is not recoverable from the collapsed residue.
\end{proposition}

\begin{proof}
A section would assign to each equivalence class a canonical representative
history. But multiple distinct histories may produce the same residue.
No choice among them is canonical without information absent from the residue.
\end{proof}

\begin{example}[Arithmetic Irreversibility]
$(2+3)$ and $(1+4)$ both produce residue $5$. From $5$ alone, neither history
is recoverable. The history is lost; the result is authoritative. This is
the Spherepop instance of the central thesis: collapse preserves embedding
(both histories map to $5$) but not recovery (the residue does not retract
onto either history).
\end{example}

\begin{theorem}[Monotonicity Under Pure Pop Sequences]
\label{thm:pop-monotone}
A sequence of admissible Pop operations with $w_i \geq 0$ satisfies
$L(\Sigma_{k+1}) \leq L(\Sigma_k)$.
\end{theorem}

\begin{proof}
Each admissible Pop removes bubble $B_i$ contributing $w_i U(B_i) > 0$.
The admissibility condition ensures all descendants have $U = 0$, so no
new load is introduced. Hence $L$ strictly decreases.
\end{proof}

\begin{remark}[Spherepop predicts its own epistemic harm failure]
Theorem~\ref{thm:pop-monotone} holds for pure Pop sequences. Meld and
Reframe may temporarily increase $L(\Sigma)$. This is the Spherepop instance
of the quantum Bayes phenomenon: the simplest operations (Pop, classical
conditioning) are monotone; richer operations (Meld, Reframe, quantum
recovery) may violate monotonicity while preserving global convergence.
The calculus predicts the pattern it is modeling.
\end{remark}

\begin{interlude}{Continuation Admissible}
The fundamental question of the Spherepop calculus is not: what does this
bubble mean? It is: does an admissible continuation exist?

The six-stage pattern of this book is organized around the same question
at every scale. The governing operations of a theory ask, at each step,
whether the current state admits a valid successor. Refuse is the answer
``yes, but not this operation.'' Collapse is the answer ``no, and none
will do.''

This is why the framework is fundamentally different from accounts that
emphasize complexity or abstraction as drivers of ontological change.
Complexity does not force enlargement. The absence of admissible continuation
forces enlargement. The Simons cone does not make minimal surface theory more
complex. It makes admissible continuation impossible within the smooth
category. The collapse is not a failure of understanding. It is a failure
of continuation.

The new ontology is not a more complex version of the old one. It is the
minimal structure in which admissible continuation is restored.
\end{interlude}

\section{Spherepop as Formal Model}

\begin{center}
\renewcommand{\arraystretch}{1.3}
\begin{tabular}{lll}
\toprule
\textbf{Stage} & \textbf{Formal Condition} & \textbf{Spherepop Realization} \\
\midrule
Refuse & $\Inadmiss(B,F) \geq \theta$, $\Admiss \neq \varnothing$
  & Pop with open descendants \\
Collapse & $\Admiss = \varnothing$
  & Collapse cycle exists \\
Enlargement & $\Onew \supset \iota(\Oold)$
  & Quotient category \\
Obstruction & $w \in \Onew \setminus \iota(\Oold)$, $w \in \F'(\Onew)$
  & Minimal obstruction subgraph \\
Accessibility & Feasible $P \mapsto w$
  & Cascade threshold $L > L_c$ \\
Commitment & Enlarged ontology accepted
  & Quotient semantics adopted \\
\bottomrule
\end{tabular}
\end{center}

The cascade threshold $L_c$ is the Spherepop control parameter: the analogue
of dimension eight in GMT and $\eta_c$ in the quantum Bayes conjecture.

% =============================================================
\part{Diagnostic Application: Foundation Models}

% =============================================================
\chapter{The Token Ontology as an Invariant System}

\epigraph{The history of artificial intelligence may be interpreted as a
search for the minimal enlarged ontology required to restore closure
under the dynamics of large-scale learning.}{}

\section{The Token Ontology in the Framework}

The token ontology expressed precisely:

\emph{Original ontology $\Oold$}: finite token sequences over a fixed vocabulary.

\emph{Governing operation $F$}: next-token prediction.
The model learns $p_\theta(x_{n+1} \mid x_1, \ldots, x_n)$, and all
capabilities are understood as consequences of this distribution.

\emph{Invariant system $\mathcal{V}$}:

$V_1$ (\emph{Representation Sufficiency}): all relevant structure can be
represented as token statistics.

$V_2$ (\emph{Predictive Sufficiency}): next-token prediction is sufficient
to explain capability formation.

$V_3$ (\emph{Behavioral Sufficiency}): deployment behavior is recoverable
from training behavior within the token-distribution description.

$V_4$ (\emph{Interpretive Sufficiency}): internal representations require
no ontology beyond token relationships to characterize model behavior
and failure modes.

All four are jointly satisfiable for small models on narrow tasks. The
question is whether they remain jointly satisfiable as models scale and
tasks generalize.

\section{Candidate Enlarged Ontologies}

Large-scale models produce structures that behave as elements of
$\Onew \setminus \iota(\Oold)$.

The \emph{embedding ontology} takes vectors in representation space as
fundamental. Geometry in that space, not token statistics, is the governing
structure.

The \emph{latent world-model ontology} takes the model's internal
representations as encoding a compressed generative model of the world.
The governing operation is inference within that world model.

The \emph{mechanistic ontology} takes specific circuits, attention patterns,
and feature directions as fundamental. These are identified by interpretability
research and may have no natural description as token statistics.

The \emph{agent ontology} takes goals and beliefs as fundamental, with the
governing operation being action selection.

Each embeds the token ontology: token prediction is a degenerate process in
each enlarged framework. Whether any of the four invariants has become
globally unsatisfiable is the open question.

\section{Candidate Obstructions}

A formal obstruction would be a capability or failure mode that cannot be
expressed while preserving all four invariants simultaneously.

\emph{Systematic out-of-distribution generalization} is a candidate
obstruction to $V_1$ and $V_2$: behaviors on inputs outside the training
distribution that the token ontology cannot predict or explain.

\emph{Alignment failures} are candidate obstructions to $V_3$: systematic
divergences between training and deployment behavior that require a different
vocabulary to characterize.

\emph{Mechanistic features with no token-level description} are candidate
obstructions to $V_4$: circuits or feature directions that cannot be
characterized as functions of token statistics.

None is yet a formal obstruction in the sense of
Theorem~\ref{thm:obstruction-nonrecovery}: no explicit construction has
proven that preserving all four invariants simultaneously is impossible for
any specific capability or failure mode.

\section{Accessibility and Phenomenology}

The accessibility structure differs qualitatively from the GMT and quantum
Bayes cases. In GMT, the obstruction was mathematically constructible. In
quantum Bayes, it may require computational exploration of large Hilbert
spaces. In the foundation model case, the obstruction region may be
\emph{empirically} inaccessible: the relevant phenomena occur in deployment
contexts that are difficult to study systematically.

This predicts a qualitatively different phenomenology: ontological commitment
in this domain will be driven by high-profile deployment failures functioning
as accessible obstructions, not by formal proofs. The engineering community
will commit to an enlarged ontology before the mathematical community has
established that collapse has occurred.

\section{Framework Predictions}

If the token ontology has collapsed, the framework predicts: the collapse is
organized by a control parameter (scale, task complexity, or both); the
transition exhibits regime structure rather than smooth degradation; the
minimal enlarged ontology is whichever candidate framework is the smallest
enlargement restoring closure; and the regularity theory of the enlarged
ontology characterizes the conditions under which token-level descriptions
remain adequate.

These are empirical predictions with specific content. They are not
predictions that the token ontology will fail, but predictions about
\emph{how} a failure, if present, will be structured.

% =============================================================
\part{General Theory and Conclusions}

% =============================================================
\chapter{A General Theory of Ontological Enlargement}

\epigraph{New objects matter because they reveal that old ontologies
were not closed under their own governing dynamics.}{}

\section{The Central Results}

\begin{metatheorem}[Closure-Recovery Principle]
\label{meta:CRP-final}
Let $(\Oold, \F)$ be a dynamical ontology. If $\cl_\F(\Oold) \neq \Oold$,
then either $\F$ must be restricted or $\Oold$ must be enlarged.
Historically, enlargement is the repeated outcome.
\end{metatheorem}

\begin{theorem}[Irreversibility of Enlargement]
\label{thm:irrev-final}
Suppose $\iota : \Oold \hookrightarrow \Onew$, recovery fails, and
$\cl_{\F'}(\iota(\Oold)) = \Onew$. Then $\iota(\Oold)$ is not dynamically
closed under $\F'$.
\end{theorem}

\begin{theorem}[Collapse as Failure of Reflection]
Ontological collapse occurs when no reflector $R : \bOn \to \bO$ compatible
with the governing operations exists. There is no canonical way to project
outputs of $\F'$ back into $\bO$.
\end{theorem}

\begin{theorem}[Regime Theorem]
Let $T$ be a universal theorem in $(\Oold, \mathcal{V})$. After enlargement,
$T$ becomes a regime theorem: valid within a recoverable subregime
$\mathcal{R} \subseteq \Onew$ defined by additional conditions, and potentially
failing outside $\mathcal{R}$. Characterizing $\mathcal{R}$ precisely is the
primary task of post-collapse regularity theory.
\end{theorem}

\begin{theorem}[Delay is Structural]
The gap between Stage~2 and Stage~6 is bounded below by the time required
to identify an accessible obstruction. This bound is governed by the structure
of the obstruction region relative to available tools, not by the sociology
of the field.
\end{theorem}

\section{Comparative Table}

\begin{center}
\renewcommand{\arraystretch}{1.3}
\begin{tabular}{p{1.6cm}p{2.6cm}p{2.6cm}p{2.6cm}p{2.6cm}}
\toprule
& \textbf{GMT} & \textbf{Quantum Bayes} & \textbf{Spherepop}
  & \textbf{Foundation Models} \\
\midrule
$\Oold$ & Smooth minimal hypersurfaces & Classical prob.\ distrib.
  & Well-nested stacks & Token sequences \\
$\F$ & Area minimization & Bayesian conditioning
  & Pop, Meld, Reframe & Next-token prediction \\
$\mathcal{V}$ & Smooth $+$ min.\ $+$ global
  & Joint $+$ unique $+$ monotone
  & Well-nested $+$ resolution $+$ gluing
  & $V_1$--$V_4$ \\
Collapse & Dim.\ 8 (established)
  & Noncomm.\ (prospective)
  & Cycle (formal)
  & Scale (diagnostic) \\
$\Onew$ & Rectifiable currents
  & Petz maps $+$ QIG
  & Quotient category
  & Embedding / world-model / agent \\
Obstruction & Simons cone
  & Petz violation (cand.)
  & Minimal cycle
  & Deployment failure (cand.) \\
Control param. & Dim.\ $n$
  & Noncomm.\ $\eta$
  & Load $L_c$
  & Model scale \\
Regime theorem & Reg.\ (dim~$\leq 7$)
  & Asympt.\ frequentism (conj.)
  & Monotone Pop
  & Token sufficiency \\
Status & Completed & Prospective & Formal model & Diagnostic \\
\bottomrule
\end{tabular}
\end{center}

\section{Open Problems}

\emph{The quantum Bayes obstruction.} Does Conjecture~\ref{conj:threshold}
hold? A simulation campaign varying Hilbert-space dimension, channel rank,
and commutator norms, testing for sharp threshold behavior in the monotonicity
violation rate, would directly address this. A positive result would establish
quantum Bayes as a second completed instance of the six-stage pattern.

\emph{Spherepop obstruction classification.} Classify all minimal obstructions
for the standard Spherepop invariant system. Determine whether minimal
obstructions have a forbidden-subgraph characterization. Develop regularity
theory for the quotient category: under what conditions does the collapsed
residue retain sufficient structure for subsequent operations?

\emph{New domains.} Apply the diagnostic to: the passage from spaces to
$\infty$-groupoids in homotopy theory; from effective field theories to UV
completions in physics; from classical logic to realizability toposes in
foundations. In each case, identify the invariant system, the governing
operations, and the obstruction structure.

\emph{Quantitative accessibility theory.} Develop a formal theory measuring
the ``distance'' from the working regime to the obstruction region. Explain
why the Simons cone was accessible within a decade while the Petz obstruction
(if it exists) has not been formalized in decades. The answer presumably
involves the difference between mathematical constructibility (GMT) and
computational explorability (quantum Bayes).

% =============================================================
\chapter{Conclusions}

\section{The Central Claim, Evaluated}

The GMT case: all six stages are documented. The invariant system was
unsatisfiable in dimension eight. GMT was available. The Simons cone was
accessible and provided the witness. The field committed. Regularity theory
characterized the recoverable regime. The pattern holds in full.

The quantum Bayes case: Stages~1 through~4 are plausible. Stage~5 is absent.
Stage~6 has not occurred. The framework correctly identifies what is missing
and specifies what a definitive result would require.

The Spherepop case: the pattern is formally instantiated in a calculus
designed to exhibit all six stages explicitly. It provides a formal model
for the structural claims, not empirical evidence but a proof of conceptual
coherence.

The foundation model case: the framework generates diagnostic questions
with specific empirical content rather than resolved answers. This is the
appropriate contribution.

\section{Structure Without Blueprint}

The framework developed here is fundamentally anti-teleological. Neither
geometric measure theory nor Bayesian inference nor the Spherepop calculus
requires a blueprint. Area-minimizing currents do not know the surface they
are attempting to become. Bayesian posteriors do not know the parameter they
will converge toward. Scope graphs do not contain their future resolutions.

Global structure emerges from local admissibility conditions rather than from
pre-specified endpoints. Ontological enlargement is not the discovery of a
hidden plan. It is the discovery that the admissibility conditions were richer
than previously recognized. The Simons cone does not reveal a deeper blueprint
for eight-dimensional geometry. It reveals that the admissibility conditions
of area minimization support more structures than the smooth category can
contain.

This is the mathematical content of the observation that the universe does
not run on blueprints. Stable structure, wherever it appears, is the residue
of local admissibility conditions operating without predetermined targets.

\section{Obstructions Do Not Create Enlargement}

Mathematical progress is often described as the discovery of new objects.
The framework developed here suggests a different picture.

New objects matter because they reveal that old ontologies were not closed
under their own governing dynamics. The Simons cone did not discover a new
kind of geometry. It revealed that smooth geometry was already insufficient
for the questions area minimization had always been asking. The Petz map does
not discover a new kind of probability. It reveals that classical probability
was already insufficient for the inference problems that quantum mechanics
poses. A collapse cycle in a scope graph does not create an irreducible
dependency. It reveals that the dependency structure was already irresolvable
within the current containment hierarchy.

Obstructions do not create enlargement. They make enlargement unavoidable.

This is why a single object can reorganize a field. Not because it is exotic
or complicated, but because it certifies, by its existence, that the field
was already living in a larger world than it knew.

We do not fade away. We resonate, as stable configurations within dynamics
that were always richer than our descriptions of them. The old ontologies
do not disappear. They persist as regimes, embedded within the larger
structures that contain them, resonating at the frequencies where the
governing operations remain within the recoverable subclass.

Return is possible. Not return to the old ontology as a universal. Return
to the conditions under which it remains valid. That is what regularity
theory is. That is what decoherence theory is. That is what scope completion
is. The enlarged framework is valuable precisely because it tells us when
we can go home.

% =============================================================
\backmatter

\chapter*{References}
\addcontentsline{toc}{chapter}{References}
\printbibliography[heading=none]

\end{document}
